\section{Replayability Grades as Evidence Licenses}
\label{sec:grades}

The activegraph-bridge project introduced five replayability levels from
observed through native.  The names arose from implementation failure modes.
We make their meaning explicit by assigning each grade a set of licensed
questions (Table~\ref{tab:grades}).

\begin{table}[ht]
\centering
\small
\begin{tabularx}{\textwidth}{@{}>{\raggedright\arraybackslash}p{0.16\textwidth}>{\raggedright\arraybackslash}p{0.32\textwidth}Y@{}}
\toprule
Grade & Required evidence & Licensed use \\
\midrule
Observed & Some inspectable trace or result & Inspection and available lineage only \\
Envelope & Faithful invocation envelope and completion, no lossy-envelope fact & Recorded-output playback; fork before invocation \\
Boundary & Envelope; mediated effects; clean reconstruction; complete known effect boundaries; no blockers & Re-execution using recorded effects; forks at effect boundaries \\
Checkpointed & Boundary plus recorded checkpoint & Resume from a restorable snapshot \\
Native & Checkpointed plus native-runtime evidence & Native replay, fork, and structural diff \\
\bottomrule
\end{tabularx}
\caption{Replay grades and their licenses.}
\label{tab:grades}
\end{table}

Let $Q(g)$ be the set of questions licensed by grade $g$.  The current taxonomy
is defined as the chain

\begin{equation}
 Q(Observed) \subset Q(Envelope) \subset Q(Boundary)
 \subset Q(Checkpointed) \subset Q(Native).
\end{equation}

This chain is a definitional summary, not a substantive theorem.  The artifact
checks its ordering mechanics, but the paper's claim is narrower: a grade
summarizes the questions supported by retained evidence, while a cut-level
assessment decides one proposed fork.  A richer taxonomy could contain
incomparable evidence sets.

\subsection{Computation from the log}

The grade is computed from projected evidence and effect completeness.  It is
not accepted as a source label.  In particular, a \code{NativeRuntime} marker
alone does not grant Native: the log must also satisfy checkpoint, boundary,
and envelope prerequisites.  Unknown footprint, uncaptured results, incomplete
or unknown lifecycles, integrity faults, lossy envelopes, unmediated effects,
and explicit hazards are blockers.

The adapter follows the same fail-closed rule.  It grants evidence only for
exact recognized source event types.  An event whose name merely contains
\code{verification} remains an unclassified domain signal.  A verification
event must contain an explicit Boolean verdict.  The bridge-specific Boundary
derivation additionally requires an effects-served count and an explicit null
divergence: the source payload's \code{divergence} field must be present and its
JSON value must be \code{null}, asserting that the bridge verifier detected no
unresolved reconstruction divergence.  This is logged evidence, not an
independent proof of equivalence.  These requirements prevent an event-name
coincidence from raising the grade.

\subsection{Grades are not monotone in time}

It is natural to view grades as achievements that only rise.  That view is
incorrect for evidence predicates.

\begin{counterexample}[Hazard downgrade]
A prefix records a faithful envelope, completion, mediated boundary, clean
reconstruction, and successful verification.  It grades Boundary.  Appending a
\code{EvidenceRecorded(HazardDetected detail)} event adds an explicit blocker,
so the extended log grades Envelope.
\end{counterexample}

The grade of a \emph{fixed log or prefix} licenses questions about that evidence.
It need not be monotone under extension.  This matters for forks: the grade of
the shared prefix can be higher or lower than the grade of the completed source
run, and a later hazard must not retroactively be ignored.

\subsection{Connection to fork safety}

The grade is a compact mechanism for checking some fork preconditions, but it
does not replace property-relative assessment.  Boundary evidence suggests
that effect results and reconstruction are available, licensing forks at known
effect boundaries.  A specific cut can still split a request from its outcome.
Likewise, a Boundary run can contain a committed OneShot effect that makes a
CounterfactualWorld claim conditional or unsound.

The validation therefore reports both the run-level grade and all cut-level
verdicts.  A grade answers ``what kinds of questions may this evidence support
in principle?''  A fork assessment answers ``does this particular cut satisfy
the conditions for property $P$?''
